\section{Returning the increment: comparison and passive layers}\label{ap:ascent}

We now return from a productive terminal path to its ancestors, in the
order $D,D-1,\ldots,1$.  The absolute increment has supplied a target
$\lambda\leq1$ separated from the desired final target by a fixed
proportional margin.  At boundedly many places in the ascent we decrease
$\lambda$ by a factor $1-\varepsilon$.  The total decrease is chosen
smaller than this margin.  If the absolute rule instead gave
impossibility, there is nothing to return.

Two features of the return will be maintained exactly.  First, the
determining slots continue to use the original prepared polynomials;
they are not replaced by approximations in their integer values.
Second, we transport the positive and negative payoff together.  A
favorable choice for one of them alone need not be favorable for their
difference.

\subsection{Paired cutoffs and the inductive score}

For $c\in\R^d$ and $r>0$, choose product cutoffs
$U^-_{r,c},U^+_{r,c}:\R^d\to[0,1]$.  The small cutoff is one when
$\|z-c\|_\infty\leq(1-\gamma)r$ and vanishes when
$\|z-c\|_\infty\geq r$.  The large cutoff is one when
$\|z-c\|_\infty\leq(1+\gamma)r$ and vanishes when
$\|z-c\|_\infty\geq(1+2\gamma)r$.  Piecewise linear transitions,
or smooth transitions with the same bounds, suffice.  Their full
coverage volumes
\[
 V_\pm=\int_{\R^d}U^\pm_{r,c}(z)\,dc
\]
are independent of $z$, and
\begin{equation}\label{ap:ascent:coverage-ratio}
 1\leq \frac{V_+}{V_-}\leq(1+C\gamma)^d,
 \qquad
 \log(2+\Lip U^\pm_{r,c})
 \ll \log(2/r)+\log(2/\gamma)+\log(2+d).
\end{equation}
Small cutoffs lie below the corresponding strict box indicator; large
cutoffs dominate the corresponding enlarged box indicator.

\begin{lemma}[Pairing the absolute patch]\label{ap:ascent:paired-cutoffs}
The absolute patch on a productive terminal path may be replaced by a
pair of product cutoffs as above, after a small proportional decrease
of its target.  The resulting score, counted with full parameter-box
mass, is positive with inverse logarithm polynomial in $p$ and the
dimension bounds.  The absolute radii have the same type of inverse
logarithmic bound, without dependence on any determining width.
The number and weights of the absolute slots are unchanged.
\end{lemma}

\begin{proof}
Write the absolute patch as
\[
 B(t)=\sum_{a\in\Z^{d_0'}}\Phi(S_t(a)),\qquad d_0'\leq d_0,
\]
as in Definition~\ref{counting:patch}, with $\Phi$ extended by zero.
Suppose its full-mass score at target $\lambda_0$ is at least $\sigma$.
Include the explicit mass of its parameter slice in $\sigma$.
For either sign, convolution with the cutoff gives
\[
 \left|V_\pm^{-1}\int\Phi(c)U^\pm_{r,c}(z)\,dc-\Phi(z)\right|
 \leq C\Lip(\Phi)\operatorname{poly}(d_0')r.
\]
The union of the supports in this comparison is still a uniqueness
box if $r$ is sufficiently small.  Consequently the sum over integer
tuples has at most one nonzero term, even on this enlarged union.
Uniform convolution error therefore bounds the error in each tested
average, without an uncontrolled sum over tuples.

Choose a slightly smaller target $\lambda$, and then choose $\gamma$
so that $\lambda V_+/V_-\leq\lambda_0$.  Finally choose $r$ so that
the displayed uniform error is smaller than a sufficiently small
multiple of $\sigma$.  Averaging the paired score over $c$ with the
nonnegative weight $\Phi(c)$ gives at least $V_-\sigma/2$.
The center region has bounded, or dimension-exponential, volume.
Hence one center has a positive paired score with the asserted
quantitative bound.  The logarithms of $V_-^{-1}$, the inverse score,
and $r^{-1}$ have polynomial bounds in the absolute patch data.
These data depend on $p$ and dimensions, not on old widths.
The chosen bounded center can either be retained in the cutoffs or
absorbed into the target polynomials.  Extending a polynomial from a
reindexed affine slice by inverse affine substitution causes no
coefficient restriction, since polynomial coefficients are unrestricted.
\end{proof}

Fix a productive prefix at layer $j$.  Its parent is an integer box
with variable $u$, and all previous path maps and pulled-back input
functions are now fixed.  In particular $C_j$ is an ordinary
polynomial on this parent.  Earlier substitutions have also been
made in the later $C_i$; we keep the notation $C_i$ for those exact
polynomials.  Fix a feasible center at layer $j$ and write
\[
 C_j(\psi(t))-c=y(\psi(t))+\beta(\psi(t)),
 \qquad b_j=\beta+b_{j,0}.
\]
The path ensemble and this identity are those of
Lemma~\ref{ap:path:ensemble}.

\begin{definition}[Ascent format]\label{ap:ascent:format}
At entry to layer $j$, the following data and properties constitute
the ascent format.
\begin{enumerate}[label=\textup{(\roman*)}]
\item A warm-cost slice of the parameter box carries a full-mass
score at least $\delta$, with $\log(1/\delta)$ warm.  Its positive
payoff contains the fixed pulled-back input $f_0(\psi(t))$, where
$0\leq f_0\leq1$.  When $j\leq s$, it also contains the current
factor $1_{\mathrm{no\ bad}}(t)$.  Its negative payoff has multiplier
$-\lambda$.

\item Both payoffs use, respectively, the small and large cutoffs of
the absolute slots and of all determining slots $b_i$ with $i>j$.
For a determining block these are cutoffs of the \emph{actual}
residual
\[
                  b_i-C_i(u,b_{\leq i-1})
\]
at a bounded constant center $l_i'$ and a radius whose logarithmic
loss relative to $w_i$ is warm at layer $i$.  Their enlarged supports
lie inside fixed padded versions of the old boxes.  The block may
also be restricted to a residue class of warm cost at layer $i$.
Payoffs are summed over their allowed integral tuples.  Each unsigned
sum, with or without $f_0$, is bounded by one.

\item Determining polynomials never use absolute slots.  Absolute
polynomials may use determining slots $b_i$ with $j<i\leq s$, but
not passive slots $i>s$.  The combined system has a fixed joint
weighted-degree bound, with determining weights exactly $i$.
Absolute weights are unchanged at passive steps and may be inflated
at active steps.  Once active determining slots are present, absolute
weights exceed their weights.

\item Polynomial extensions and determining identities are formal
identities.  Equality of integer slot arguments is never inferred
merely from closeness of the corresponding real polynomials.
\end{enumerate}
\end{definition}

The terminal absolute return, after
Lemma~\ref{ap:ascent:paired-cutoffs}, supplies the initial instance.
The surviving paths of Proposition~\ref{ap:scalar:forward} are
modular-good at every layer.  All productive subsets used below are
subsets of these surviving paths; the active argument therefore starts
with modular-good children even when its later selections have cold cost.
The output of a layer will have this format at layer $j-1$, with
the current raw block of $d_j$ determining slots included.  Its
full-mass score may contain a factor $m_j$.  Thus its inverse
logarithm may charge $Q_j$, which is warm at the next ancestor.
The relative loss in the new width at layer $j$, by contrast, must
be warm at layer $j$ itself.

We only need the return at parents having a warm-inverse fraction
of productive children.  This follows by conditional averaging:
if the total productive probability is at least $\kappa$, a set of
parents of probability at least $\kappa/2$ has conditional productive
probability at least $\kappa/2$.  At each such parent, average once
more over its feasible centers to fix a center with a comparable
conditional probability.  Repeat at the bounded number of layers.
Inside one fixed parent argument a subsequent choice may lose a
cold-pre fraction of its paths; this does not multiply the already
obtained measure of parents where an output exists.

These choices require no simultaneous measurable choice of a
polynomial witness over a continuum of centers.  For fixed input
data, earlier integer affine maps, integral lift polynomials and
flags have finite possibilities.  They index the discrete states
on which the rank tests and descendant witnesses depend; center
integration itself is measurable.  After fixing a center, the path
law again has finite discrete support.  Witnesses can consequently
be selected at these states after the probability estimates.

\subsection{The comparison measure}

The score in Definition~\ref{ap:ascent:format} is an average on a
child parameter box and still contains polynomial data chosen
separately on each child. At a passive layer,
Proposition~\ref{ap:passive:replacement} replaces these data by one
ambient polynomial; scalar comparison then transfers site payoffs
from one finite list. At active layers, Section~\ref{ap:active}
uses this comparison at the degree-zero end of its induction.
We first define the weighted parent average needed for these steps.
The inactive
parameters retain their actual integer values.  Each parallel
parameter has a real copy and a residue copy; these are averaged
independently before the site and incidence constraints are imposed.
The common site then relates the two copies.

For a modular-good path and each fixed actual $t_I$, let
$d(b;t_I)$ be the density of $b(t)\bmod\mathfrak M$, relative to uniform
output measure, when $t_\parallel\bmod\mathfrak M$ is uniform.
Lemma~\ref{cube:modular-decay} gives a warm cap for $d$, and an
absolutely convergent Fourier expansion whose tail above individual
character order $T$ is bounded by a warm constant times an arbitrarily
large fixed inverse power of $T$. A pre-cost stride restriction on
$t_\parallel$ has the same conclusions with a pre-cost prefactor.
Different strides can be combined into one divisible stride and
disjoint residue classes. These conclusions also hold for any
complementary subset of the primes. In particular none of these caps
grows with $Z$.

Fix a modular-good path. Define the nonnegative functional $\mathcal K$
on bounded site-and-label functions by the following integration rule.
\begin{enumerate}
\item Use the full-parent average $\E_{u\in U}$ with the indicator of
the buffered uniqueness chart, not the average conditioned on that
chart. Independently average over one actual
integer vector $t_I$ in its original ranges, and over an independent
uniform $t_\parallel^{\rm res}\bmod\mathfrak M$. Impose
\begin{equation}\label{ap:scalar:incidence}
 z_I(X)=z_I(t_I),\qquad
 b(X)=b(t_\parallel^{\rm res},t_I)\pmod{\mathfrak M},
\end{equation}
and multiply by
\begin{equation}\label{ap:scalar:normalization}
                   m_j^{-1}\Bigl(\prod_{i\in I}k_i\Bigr)
                                  \mathfrak M^{m_b}.
\end{equation}
\item Use normalized real parallel parameters in $[0,1]$. Keep the
spatial kernel output and the continuous and active principal lift
outputs, with the path's actual constant shifts. Their required output is
\begin{equation}\label{ap:scalar:real-output}
                (u/H,y_W(u)/v,z_{\rm act}(u)/k).
\end{equation}
Choose one spatial kernel block and one variable in a dedicated
product block for each lift output direction as pivots. Integrate
the remaining real parameters, solve linearly for the pivots giving
\eqref{ap:scalar:real-output}, and retain only solutions in their
parameter ranges. Weight each solution by the inverse absolute
determinant; a zero determinant contributes nothing. Evaluate the
real labels at this solved parameter vector. This defines the
integration on each output fiber even for an arbitrary bounded
function of the labels.

Use the orthonormal coordinates and measure convention defining $m_j$.
Nonkernel spatial displacement and nonprincipal lift perturbations
are omitted only in this real density, not in the integer congruences.
\item A slice imposes its real interval and residue conditions on the
two parallel copies, and its actual integer membership condition on
$t_I$. Unless explicitly normalized, its mass is retained. Extra
normalized labels use the real parallel copy and $t_I/T_I$; extra
integer labels use the residue parallel copy and the same actual
$t_I$.
\end{enumerate}

\begin{lemma}[Comparison measure and forecasting weights]
\label{ap:scalar:measure}
For site functions alone these are the normalized-chart forecasting
weights of Equation~\eqref{cube:forecast}. The unrestricted site
weight is at most $B m_j^{-1}$ for a warm $B$. Normalized pre-cost
slices have cap $B_{\rm pre}m_j^{-1}$ and uniform approximations, on
buffered charts, by pre-cost combinations of order-zero site twists
times $m_j^{-1}$. Both approximation error and coefficient bounds
are measured in these normalized units and are independent of $Z$.
Smoothness is measured in $(u/H,y/v)$.
\end{lemma}

\begin{proof}
For site functions, the pivot integral is the output density by the
change-of-variables formula. Spare
independent smoothing blocks ensure that this formula gives the same
continuous version of the density at every output, also with
rectangular interval restrictions: a possibly singular block law is
convolved with a bounded Lipschitz density. Thus no value of an
arbitrary almost-everywhere density representative is being evaluated
at a true site. Lemma~\ref{cube:product-block} supplies the required
smoothing, while independent nonsingular spatial kernel blocks give
integrable Fourier decay with a first moment.
This continuity assertion concerns the site weight. Arbitrary label
insertions are evaluated by the explicit pivot integral above.

The same convolution observation controls deletion of pivot choices
with very small determinant, or pivot coordinates in narrow boundary
strips. Condition on the exceptional pivot group before convolving
with spare blocks. Its contribution to every output density is bounded
by a controlled constant times its parameter mass. The determinant is
a product of bounded-degree factors and fixed nondegenerate principal
amplitudes; keeping its factors outside appropriately small strips
therefore has polynomial logarithmic cost. These facts concern density
in normalized output coordinates, not a point mass at a critical
image.

We now check the cap and approximation. Conditional on $t_I$, the
congruence density has the cap following
Equation~\eqref{ap:scalar:modular-goodness}. The inactive product-block
law has a bounded density relative to its grid volume. Its extra
linear column is handled by conditioning and translating; omitted
axes have warm-bounded $k_i$ and are enumerated. Multiplying these
caps by Equation~\eqref{ap:scalar:normalization}, and by the bounded
real-output density, gives $B m_j^{-1}$. The corresponding slice
proof charges its explicit normalization and pre-cost restrictions.

For the approximation, truncate the absolutely convergent congruence
Fourier series by individual character order. A term of order $q$
has coefficient depending only on $t_I\bmod q$, because $b$ has
integer coefficients. Split $t_I$ into these residue classes. The
order-zero product-block expansion on each class gives the remaining
smooth and residue twists; an axis too short after the restrictions
has pre-bounded $k_i$ and can instead be enumerated. Characters of
sheet and deck coordinates are characters of $\beta$ at the same
modulus. The real kernel and dedicated blocks have the stated
normalized smoothness. Summing the finitely many retained orders and
dividing the error among them preserves a pre bound. This is the
forecasting-weight proof preceding
Proposition~\ref{cube:conditioned-transfer}, now in units $v,k_i$.
Its absolute Fourier tail and all caps are uniform in $Z$.
\end{proof}

\subsection{Proof of the selected-plane estimates}

The comparison measure retains the congruence data of the whole path.
The identity in Equation~\eqref{ap:scalar:plane} provides a family of
sections of that data: the variable $s$ changes the output by a
prescribed amount, while $h$ changes the inactive linear variables.
The selection tests ensure both controlled counting weight and the
translation invariance needed when these sections are used in extraction.

\begin{proof}[Proof of Lemma~\ref{ap:scalar:flags}]
We first prove the finite-field assertions, including the characteristic
dependence. Temporarily write $P=\varpi$, $N=v_f$, and fix the other
coefficients. Every nonconstant pure-free coefficient of every output
is independently uniform. An affine change of free variables acts
invertibly on polynomial coefficients modulo constants, in any
characteristic. For a full-rank $V$, restriction to $x+V(s,h)$ maps
these coefficients surjectively onto all polynomials in $(s,h)$ modulo
constants of the indicated degree. This follows, for example, by
choosing a linear left inverse of $V$. Consequently
Equation~\eqref{ap:scalar:plane} has probability exactly $P^{-J}$.
The fixed constant coefficients do not affect this assertion, since
the identity subtracts $b(x,r_0)$.

Let $N_0$ be the unselected solution count at a given basepoint and
$H=P^{N\ell-J}$. Restrictions to two direction spaces with direct sum
have independent nonconstant coefficients. The proportion of
rank-deficient matrices is $O(\ell P^{\ell-N})$, and the proportion
of pairs whose column spaces do not have direct sum is
$O(\ell P^{2\ell-N})$. Using the trivial joint probability bound on
these bad pairs gives
\begin{align}
 \E(N_0/H)&=1+O(\ell P^{\ell-N}),\label{ap:scalar:plane-first}\\
 \E\bigl[(N_0/H)^2\bigr]&=1+O(\ell P^{2\ell-N+2J}).
                       \label{ap:scalar:plane-second}
\end{align}
The harmless error in the first line is also absorbed into the second.

Condition next on Equation~\eqref{ap:scalar:plane} for a specified
full-rank first matrix. At a specified translated point of its plane,
two further direction spaces retain independent uniform nonconstant
restrictions whenever all three free direction spaces have direct sum.
Indeed use basis coordinates for their direct sum, centered at that
translated point; translation along the first plane changes its fixed
restriction but leaves the other nonconstant coefficient arrays free.
The same first- and second-moment argument has error at most
\begin{equation}\label{ap:scalar:conditional-moment}
                   O(\ell P^{3\ell-N+2J}).
\end{equation}
This bound allows the crude estimate on all nondirect triples.

Here is the Jacobian check. In free coordinates adapted to the first
plane, write the transverse variables as $y_1,\ldots,y_{N-\ell}$.
Conditional on its restriction, the transverse derivatives in these
directions are independent uniform vectors in
\[
          E_{\rm der}=\bigoplus_{i=1}^{m_b}
                      \operatorname{Pol}_{\le d_i-1}(\mathbb F_P^\ell),
             \qquad \dim E_{\rm der}=n_{\rm der}.
\]
The probability that $N-\ell$ such vectors fail to span is at most
$O(P^{n_{\rm der}-(N-\ell)})$, by a union bound over nonzero dual
vectors. If they span, varying the $q$th column of $V$ realizes that
space multiplied by the $q$th plane coordinate. Every nonconstant
monomial of degree at most $d_i$ has such a factor, so these variations
span all $J$ coefficient equations. This proof differentiates only
terms $y_a(s,h)^\alpha$, whose transverse exponent is one. It uses
neither factorial division nor a derivative of a high power and is
valid also when $P\le j$.

There are $P^\ell$ translated basepoints on a plane. Apply Chebyshev's
inequality at tolerance $\epsilon_P/2$ using
Equation~\eqref{ap:scalar:conditional-moment}, and union bound over
these points together with the derivative failures. Since $J$, $\ell$,
and $n_{\rm der}$ do not involve $N$, taking $N$ a sufficiently large
polynomial in the structural dimensions makes the conditional
probability that a solution is unselected at most, for example,
$c_\star\epsilon_P^2P^{-20}$. Summing over solutions bounds the expected
deleted count divided by $H$ by the same quantity. Markov's inequality
for the deleted count and
Equations~\eqref{ap:scalar:plane-first}--\eqref{ap:scalar:plane-second}
for the full count prove the asserted concentration. The exact upper
test in the selection rule also gives $w_{\varpi,1}\le1+\epsilon_\varpi$:
if any matrix is selected, that test holds at the original basepoint.

The formal plane identity persists at every point of its plane.
Both tests defining selection inspect that whole plane, so the
selected incidence relation is translation invariant. At a selected
solution modulo $\varpi$, the Jacobian is surjective. A lift from
depth $a$ to depth $a+1$ requires solving $J$ independent linear
equations in the $N\ell$ next digits of $V$. There are exactly
$\varpi^{N\ell-J}$ solutions. Induction proves the claimed depth
independence at every compatible lift of the basepoint, regardless
of its higher digits or of those of the fixed polynomial coefficients.

The product bounds follow from
$\sum_\varpi\epsilon_\varpi=O(c_\star)$. Chinese remaindering makes
the nonconstant coefficient arrays at different primes independent,
even after conditioning on the remaining nonconstant data. This gives
the joint flag bounds. Notice that all equations and their Jacobians
ignore the constant coefficient of $b$.

For the actual path law, Lemma~\ref{ap:path:ensemble} gives the needed
finite-residue comparisons. Uniformity is required in all nonconstant
pure-free degrees, not only the leading degree. On moderate axes these
are precisely the added pure-free perturbations; the other axes use
the original continuous/enormous laws and uniform deck choices. The
late range choices make every required coefficient interval broad at
each individual modulus.

For completeness, the proportional estimate with an arbitrary site
insertion needs the analytic-removal assertion as well as the residue
law. Fix a parameter value $t$ and average the unnormalized tilted
one-site integral over all centers. Its normalizer is $1+o(1)$
uniformly. Center averaging removes the constant-coefficient density.
Enumerate the spatial nonconstant residues of the flag event and
Fourier approximate its nonconstant lift-coefficient condition on the
individual prime cover. A nontrivial coefficient mode factoring
through the site has its site row as its constant row. Since the
flag condition has constant row zero, such a surviving factoring mode
is trivial. Lemma~\ref{sampling:analytic-removal}, in the form verified
in Lemma~\ref{ap:path:ensemble}, removes all other modes even with the
arbitrary site insertion. What remains is the Haar-chart flag
probability conditional on the specified spatial nonconstant
coefficient residues modulo $\varpi$. Denote this residue array by
$A_{\rm sp}$. Analytic removal is applied on each residue class, not
after fixing the exact values of all spatial slopes. In the remaining
unweighted spatial average, the independent root law has an upper
density arbitrarily close to uniform for every fixed slope tuple.
Hence, if $q_\varpi(A_{\rm sp})$ is the conditional Haar-chart flag
probability, the bound on each residue class, for a nonnegative
insertion $F$, is
\[
 q_\varpi(A_{\rm sp})(1+o(1))\E_U F.
\]
Only after averaging the broad spatial slope residues $A_{\rm sp}$ do we use
$\E q_\varpi(A_{\rm sp})\le c_\star\varpi^{-10}$; the conditional
probability itself need not be small. Summation over individual
primes proves the result.
The approximation error per prime can be chosen small enough that its
sum through $Z$ is still arbitrarily cold-small. No simultaneous law
modulo $\mathfrak M$ was used.
\end{proof}

\subsection{Scalar transfer with labels and flags}

A menu below is a finite list of bounded functions of the true site,
with no regularity assumed in that variable. At a fixed parent and
center, the same menu is used for every path. Its entries may depend
on ambient maps already selected in the return argument; a path may
then select an entry and a permitted slice. The logarithmic
cardinality, cap and required accuracy are forecast before the late
sampler choices. This forecasts bounds on the eventual menu, rather
than requiring the ambient maps themselves to be chosen before the
ensemble is constructed.

\begin{proposition}[Scalar comparison]\label{ap:scalar:comparison}
Suppose a set of productive modular-good paths has probability at least
$\tau$, with $\log(1/\tau)$ pre bounded, and each has a specified
pre-cost slice. Consider bounded scores depending on site and auxiliary
labels. Assume that, after partitioning into pre-count many pre-cost
slices, each score is uniformly approximable at the required pre
accuracy, both on actual path arguments and on the joint support of
$\mathcal K$, by the same menu of bounded site functions. A bounded
number of successive increases of these pre bounds is allowed.

Then, outside a set of paths smaller than the productive probability,
the actual scalar averages and their $\mathcal K$ integrals differ by
an arbitrarily prescribed pre-small error. This holds simultaneously
for a pre-count menu of scores and for path-dependent slice choices.
At active layers it also holds after multiplying a score by the
no-bad-flag indicator. It is valid either for normalized slice averages
or with full-box mass retained, paying the explicit slice fractions.
The budget is warm for passive uses and cold pre otherwise.
\end{proposition}

\begin{proof}
First omit flags. Enlarge the degree-zero local testing seminorm by
the supremum of pairings with all normalized sliced forecasting
weights of Lemma~\ref{ap:scalar:measure}. Their cap and uniform
order-zero twist approximation give detection for the enlarged
seminorm. Apply the separation argument of
Lemma~\ref{cube:separation}, using Theorem~\ref{ap:path:detection}.
In the warm case its inner product is the normalized chart measure,
and the one-site cap and tail are taken relative to that measure.
Every fixed bounded masked site input therefore has an order-zero
twist model whose residual is small both in the pathwise slice
seminorm and against every forecasting weight. Fix the number of
atoms, their coefficient mass, and their complexity before $Z$ is
chosen. Markov's inequality and a union bound over the menu make all
the actual residual errors small outside a set of paths of probability
less than any prescribed fraction of $\tau$. The supremum in the
seminorm permits slices to be chosen after the path.

Choose $Z$ to include every prime power needed by these atom moduli,
slice strides, and label periods. For a model twist
\[
          J(u)=F(u/H,y(u)/v)R(u,\beta(u)\bmod q),
\]
the actual and comparison integrals agree at the required accuracy.
Here are the normalization and modulus details. Truncate the rational
Fourier density at individual order $q'$. Each retained term and the
atom require only the cover of modulus $\operatorname{lcm}(q,q')$.
Apply single-site Haar-cover equidistribution on this cover. Exact
inactive equalities can be extended by buffered cutoffs at the actual
grid resolution; since inactive $k_i\le L^j$, this has permitted late
cost. In orthonormal lift coordinates, Haar sheet measure is Lebesgue
measure on $W$ times counting measure on the sheet coordinates, divided
by $\covol(W\cap\Z^{d_j})$. Multiplication by
$m_j^{-1}\prod_{i\in I}k_i$ leaves normalized Lebesgue measure in
$y_W/v$, exact inactive counting measure, and active grid volumes
$\prod_{i\notin I}k_i^{-1}$. The spatial and deck residues are
independent uniform variables. On each active axis the normalized
spacing in the needed residue class is at most
$\operatorname{lcm}(q,q')/k_i$. Since $k_i\ge(AL)^j$, the later choice
of $A$ makes the grid sum converge uniformly to the real density.

Restore the rational density by its uniform absolute tail. Its integral
against $R$ depends only on the marginal modulo $q$, not on the whole
$\mathfrak M$. Actual Riemann summation on each parallel parameter in
its prescribed residue class gives that same marginal jointly with
the real argument. The omitted nonkernel spatial terms and normalized
lift perturbations have size $O(\sigma)$ and therefore change the
smooth input negligibly. Every residue input throughout uses the true
integer polynomial. No rounding of a principal real output is used to
infer its residue. These are precisely the individual-cover steps of
Equation~\eqref{cube:model-comparison}, with the present normalization.
Combining them with the two model residual estimates proves the
unflagged assertion.

The common-menu hypothesis is stable under the label operations we
shall use. Independent bounded-length integer axes can be enumerated;
integer labels can be split at pre-cost moduli dividing $\mathfrak M$;
and pre-Lipschitz dependence on long normalized labels can be meshed
inside the original slice while the integer labels are held fixed.
Use the same mesh value on actual and comparison supports. Axes too
short for a mesh or stride are enumerated exactly. Parallel lengths
were chosen after these requirements, so endpoint rounding there is
negligible. These operations cost a union bound over menu entries,
not a simultaneous pigeonhole of all their pathwise choices.

We next prove the flagged assertion and record its choice order.
At a given approximation step choose a large pre threshold $T$, before
choosing $Z$, and replace the no-bad indicator by
\begin{equation}\label{ap:scalar:truncated-flags}
 \sum_{\substack{d\le T\text{ squarefree}\\
                 \varpi\mid d\Rightarrow\varpi\le Z}}
       \mu(d)\,1_{\{\text{bad at every }\varpi\mid d\}}.
\end{equation}
This is exact if the product of the bad primes is at most $T$.
Otherwise its absolute error is at most $1+T$. The exceptional event
has a witness consisting either of one bad prime $\varpi>T$, or of a
squarefree product $d\in(T,T^2]$ of bad primes: in the latter case
multiply primes at most $T$ until the product first crosses $T$.
Lemma~\ref{ap:scalar:flags} and summation give expected exceptional
mass $O(T^{-9})$ in the uniform model and hence error $O(T^{-8})$.
The same conclusion holds in actual unrestricted path/parameter
averages at fixed center, up to negligible errors, by the individual
finite-residue laws for nonconstant coefficients. For a normalized
slice of probability $\rho$, divide the full-mass error by $\rho$;
this charges only the permitted pre factor $\rho^{-1}$.

We need the same statement for $\mathcal K$, where incidence weights
cannot simply be ignored. On modular-good paths, retain the matching
at the witnessing primes and bound matching at all complementary
primes by their modular Fourier cap, conditional on $t_I$. Flags see
only first digits, but retain matching to the full tested depth at
each witness prime. For a squarefree witness $d\le T^2$, this full
modulus is
\begin{equation}\label{ap:scalar:witness-modulus}
 Q_d=\prod_{\varpi\mid d}\varpi^{a_\varpi},\qquad
 \log Q_d\le\frac{2\log T\log Z}{\log2}.
\end{equation}
For a single prime witness its full modulus has logarithm at most
$\log Z$. Thus each comparison has allowed late cost after $Z$,
even though $\mathfrak M$ itself need not be affordable. The prime
set of a witness does not depend on the site.

On the $Q_d$ cover integrate the real output density and the exact
inactive constraints with the prescribed residue data, by the
individual-cover computation above, increasing $A$ also for these
moduli. Integration removes the retained incidence-density factor
and leaves uniform parameter residues on the witness primes. The
actual $t_I$ can be held fixed throughout. Complementary parallel
residues were bounded by their separate cap, and the Chinese
remainder theorem separates them from the witness residues.
This cover calculation is uniform in the fixed path data: the real
principal amplitudes have bounded normalized ranges, exact inactive
cutoffs have grid resolution $k_i^{-1}$ with $k_i\le L^j$, and all
integer coefficients in the witness condition are reduced modulo
$Q_d$. Their Fourier costs are therefore among the forecast cold and
late costs, independently of unreduced integer coefficient sizes.
In particular, writing $E_d$ for the witness flag event, the full-mass
inequality on a modular-good path is
\begin{equation}\label{ap:scalar:witness-cancellation}
 \mathcal K(1_{E_d})
 \le C_{\rm pre}\,
       \E_{t_I}\E_{t_\parallel\bmod Q_d}1_{E_d}+o(1).
\end{equation}
There is no factor $Q_d^{m_b}$ left in this inequality: the retained
matching factor cancels against uniform output mass on the cover,
and the remaining full real density and normalized inactive-grid
mass integrate to at most one. A normalized slice instead divides
this bound by its slice probability. The errors can be chosen
uniformly small enough to sum over all witnesses, since their number
is at most $Z+T^2$ and $\log Z,\log T$ have the prescribed bounds.
Average over paths and drop their goodness restriction for this
nonnegative upper bound. The finite-residue comparison at the witness
then gives its uniform flag probability. Summing witnesses gives
$O(T^{-8})$ error, now with a fixed pre prefactor from $D'_0$ and the
normalized real densities. This is the required forecast-side bound.

Markov's inequality makes both truncation errors small on enough
productive paths, taking $T$ large relative to the current score,
probability, and cap. Each retained intersection in
Equation~\eqref{ap:scalar:truncated-flags} uses a modulus $d\le T$ and
is constant on a further pre-cost residue slice, separately for each
path. The unflagged comparison therefore applies to it. Subsequent
models and menus may charge $\log T$: the truncation just proved used
only the unexpanded score and the current cap and accuracy, not those
later menu sizes. A bounded number of required truncations and models
can consequently be forecast successively before $Z$. Then $Z$ is
chosen to contain their periods, and only then are the individual
witness comparisons and late sampler parameters fixed. This order
avoids any dependence on the unaffordable size of
$\mathfrak M$. The proof also shows that actual-path flag truncation
alone may be used in positive-degree detection and pivot selection.
\end{proof}
